% Generated draft: requires mathematical and source review.
\documentclass[11pt]{article}
\usepackage[margin=25mm]{geometry}
\usepackage{fontspec}
\setmainfont{Latin Modern Roman}
\usepackage{amsmath,amssymb,mathtools,unicode-math}
\setmathfont{Latin Modern Math}
\usepackage{enumitem,needspace,xurl,hyperref}
\setlength{\parindent}{0pt}
\setlength{\parskip}{5pt}
\setlength{\emergencystretch}{3em}
\newcommand{\N}{\mathbb{N}}
\newcommand{\Z}{\mathbb{Z}}
\newcommand{\Q}{\mathbb{Q}}
\newcommand{\R}{\mathbb{R}}
\newcommand{\C}{\mathbb{C}}
\newcommand{\F}{\mathbb{F}}
\newcommand{\A}{\mathbb{A}}
\newcommand{\PP}{\mathbb{P}}
\newcommand{\E}{\mathbb{E}}
\newcommand{\eps}{\varepsilon}
\newcommand{\dd}{\,\mathrm{d}}
\begin{document}
\section*{UnsolvedMath ID: 30005074}
\textbf{Source problem code:} catalogue entry
\subsection{Definitions and mathematical statement}
A topological quantum field theory is abbreviated TQFT. The supplied extract concerns invariants arising from semisimple TQFTs in even dimensions, but does not specify the definition of semisimplicity, the category of manifolds, the additional structures on manifolds, or the stabilization operation defining stable diffeomorphism.

The extract attributes the following assertion to its Theorem A: even-dimensional semisimple TQFTs yield stable diffeomorphism invariants. Its mathematical content can be displayed only schematically as
\[
M\sim_{\mathrm{st}}N\quad\Longrightarrow\quad I_Z(M)=I_Z(N),
\]
where \(Z\) denotes a semisimple TQFT covered by that theorem, \(I_Z\) denotes its associated manifold invariant, and \(M\sim_{\mathrm{st}}N\) denotes the stable diffeomorphism relation intended in the report. The implication is intended for every such theory and every pair of manifolds in its applicable domain. These symbols are placeholders for the missing source definitions and hypotheses; this display is not a recovered, fully specified theorem. In particular, the supplied record does not determine the allowed dimensions beyond their being even, the manifolds' boundary conditions, or the structures that the diffeomorphisms must preserve.

The extract asks, ``How strong are these invariants?'' It supplies no criterion of strength and no conjectural converse, separation property, or classification target. It additionally describes stable diffeomorphism as a much coarser relation than diffeomorphism and its classes as better understood, but the scope and supporting statements are absent. Its final sentence refers to a consequence of a classical result without supplying that consequence.

\textbf{Unresolved target.} No precise open mathematical problem can be recovered from this record. The supplied statement-verification field identifies the question as a rhetorical heading answered by the report's theorem, whereas the registry status and dated literature assessment label the item open. Neither a new open problem nor a stronger theorem is asserted here.
\subsection{Short English statement}
The extract says that semisimple TQFT invariants in even dimensions are unchanged under the report's notion of stable diffeomorphism. Essential definitions and hypotheses are missing. The question about their strength has no specified mathematical target, and the supplied verification identifies it as a rhetorical heading rather than an open problem. This statement-only draft requires human mathematical review.
\subsection{Sources}
\begin{itemize}
\item Local UnsolvedMath registry, permanent ID \(30005074\), title \emph{Strength of Semisimple TQFT Invariants}, source problem code \texttt{catalogue entry}; supplied collection: Oberwolfach Reports Open Problems. No individual proposer is supplied; the proposed year is recorded as 2022.
\item Supplied report citation: \emph{Mini-Workshop: Recent Developments in Representation Theory and Mathematical Physics (2023)}, Oberwolfach Reports, pp.\ 817--848, DOI \url{https://doi.org/10.4171/owr/2022/16}. The extracted text includes the page marker 833. The supplied literature-source label is \emph{Oberwolfach Report 2022/16}, with the same link. The citation's 2023 designation and the record's 2022 report designation are retained without reconciliation.
\item The extract attributes Theorem A to reference [3] in dimension four and reference [4] in arbitrary even dimensions. It cites [6] for background on stable diffeomorphism classes and refers to a classical result [7]. The bibliographic identities of [3], [4], [6], and [7] are not supplied.
\item Literature assessment checked at \texttt{2026-08-22}: this is supplied dated curation, not independently verified. It labels the item open, says that the extracted question lacks a precise theorem or criterion, and says the report does not resolve it. The supplied background instructs independent verification of time-sensitive status and cited claims.
\item The supplied statement-verification field instead says that the official report answers the rhetorical heading with a theorem characterizing the invariants' stable-diffeomorphism sensitivity and does not state an open problem. The supplied statement status is \texttt{unrecoverable}. This verification claim has not been independently checked here.
\end{itemize}
\end{document}
